\chapter{Representasi Grup Berhingga}\label{ch:b3:representations}

Untuk memahami grup abstrak, buatlah ia beraksi pada ruang vektor
lalu pakailah aljabar linear --- nilai eigen, trace, hasil kali dalam ---
pada aksi itu. Program ini, yang disebut \emph{teori \omterm{def:b3:representations:rep}{representasi}},
luar biasa mujarab bagi grup berhingga atas $\C$: setiap
\omterm{def:b3:representations:rep}{representasi} terpecah menjadi \omterm{def:b3:representations:rep}{representasi tak tereduksi} (Maschke),
yang \omterm{def:b3:representations:rep}{tak tereduksi} itu dipakukan oleh \emph{\omterm{def:b3:representations:character}{karakternya}}
(yakni tracenya), dan \omterm{def:b3:representations:character}{karakternya} memenuhi relasi ortogonalitas yang
membuat perhitungannya mekanis. Bab ini membangun kalkulus itu
beserta mahakarya pertamanya --- \omterm{def:b3:representations:table}{tabel karakter} grup
kecil --- dan soal akhir pekan memanen sebuah teorema yang jauh melampaui
jangkauan teori grup telanjang: teorema $p^aq^b$ Burnside.
Sepanjang bab ini $G$ adalah grup berhingga dan semua ruang vektornya
berdimensi berhingga atas $\C$.

\section{Representasi, Maschke, Schur}

\begin{definition}\label{def:b3:representations:rep}
Sebuah \emph{representasi}\index{representasi} dari $G$ adalah
morfisma $\rho \colon G \to GL(V)$ untuk sebuah ruang vektor-$\C$ $V$;
sedangkan $\dim V$ disebut \emph{derajatnya}. Subruang $W \subseteq V$ disebut
\emph{invarian} jika $\rho(g)W \subseteq W$ untuk setiap $g$;
pembatasannya menjadikan $W$ sebuah \emph{subrepresentasi}. Representasi $\rho$ disebut
\emph{tak tereduksi}\index{representasi tak tereduksi} jika $V \neq
0$ dan satu-satunya subruang invariannya adalah $0$ dan $V$. Sebuah
\emph{morfisma} antara $(\rho, V)$ dan $(\sigma, W)$ adalah pemetaan linear
$f\colon V \to W$ dengan $f\rho(g) = \sigma(g)f$ untuk setiap $g$
(keekuivarianan); dan $f$ yang bijektif disebut \emph{isomorfisma}.
\end{definition}

\begin{example}\label{ex:b3:representations:examples}
(a) Derajat $1$: yaitu morfisma $G \to \C^\times$.
(b) Adapun \emph{representasi reguler}\index{representasi
reguler}: $V = \C^G$ dengan basis $(e_h)_{h \in G}$ dan
$\rho(g)e_h = e_{gh}$; berderajat $\abs G$.
(c) Aksi permutasi $G$ pada himpunan berhingga $X$ memberikan
\emph{\omterm{def:b3:representations:rep}{representasi} permutasi} pada $\C^X$: $\rho(g)e_x =
e_{g\cdot x}$.
(d) Grup $S_n$ beraksi pada $\C^n$ dengan mempermutasikan koordinatnya; hiperbidang
$\{\sum x_i = 0\}$ bersifat invarian: itulah \emph{\omterm{def:b3:representations:rep}{representasi}
baku}, yang berderajat $n - 1$.
\end{example}

\begin{theorem}[Maschke]\label{thm:b3:representations:maschke}
Setiap subruang invarian $W$ dari \omterm{def:b3:representations:rep}{representasi} $(V, \rho)$
mempunyai komplemen yang invarian. Akibatnya setiap \omterm{def:b3:representations:rep}{representasi}
merupakan jumlah langsung \omterm{def:b3:representations:rep}{representasi tak tereduksi}
(\emph{kesemisederhanaan}).\index{teorema Maschke}
\end{theorem}

\begin{proof}
Misalkan $p \colon V \to V$ \emph{sembarang} proyeksi yang berpeta $W$
(ambil komplemen apa pun). Rata-ratakan ia atas grupnya:
\[
\tilde p = \frac1{\abs G}\sum_{g \in G} \rho(g)\,p\,\rho(g)^{-1}.
\]
Setiap sukunya memetakan $V$ ke dalam $W$ (sebab $W$ invarian), dan menetapkan $W$
titik demi titik: sebab untuk $w \in W$, $\rho(g)^{-1}w \in W$, lalu $p$ menetapkannya,
dan $\rho(g)$ membawanya kembali --- jadi $\tilde p$ kembali menjadi
proyeksi pada $W$. Ia ekuivarian: sebab untuk $h \in G$,
$\rho(h)\tilde p\rho(h)^{-1}$ hanya mengindeks ulang jumlah yang sama. Karena itu $\ker
\tilde p$ merupakan komplemen invarian bagi $W$. Mengiterasikannya pada
setiap sukunya (dimensinya berhingga) menguraikan $V$ atas \omterm{def:b3:representations:rep}{representasi tak tereduksi}.
\end{proof}

\begin{theorem}[Lema Schur]\label{thm:b3:representations:schur}
Misalkan $f \colon V \to W$ sebuah morfisma antara \omterm{def:b3:representations:rep}{representasi} \emph{\omterm{def:b3:representations:rep}{tak
tereduksi}}. Maka $f = 0$ atau $f$ merupakan isomorfisma; dan jika
$(V,\rho) = (W,\sigma)$, maka $f = \lambda\,\mathrm{id}$ untuk suatu
$\lambda \in \C$. Karena itu
$\dim\operatorname{Hom}_G(V, W)$ bernilai $1$ bila $V \cong W$, dan $0$
bila tidak.\index{lema Schur}
\end{theorem}

\begin{proof}
Kernel $\ker f$ dan peta $\operatorname{im} f$ bersifat invarian (karena keekuivarianan),
sehingga masing-masing berupa $0$ atau seluruhnya: jadi entah $f = 0$, entah $f$ injektif
dengan peta penuh. Jika $V = W$: maka $f$ mempunyai nilai eigen $\lambda$ (sebab $\C$
\omterm{def:b3:galois:closure}{tertutup secara aljabar}); lalu $f - \lambda\,\mathrm{id}$ menjadi morfisma
$V \to V$ yang tak injektif, sehingga $0$. Untuk pencacahan dimensinya
ketika $V \cong W$: tetapkan satu isomorfisma $u$, lalu sembarang morfisma
$f$ memberikan endomorfisma $u^{-1}f = \lambda\,\mathrm{id}$: jadi $f =
\lambda u$.
\end{proof}

\section{Karakter dan ortogonalitas}

\begin{definition}\label{def:b3:representations:character}
Yang disebut \emph{karakter}\index{karakter} dari $(V, \rho)$ adalah
$\chi_\rho(g) = \operatorname{tr}\rho(g)$. Ia memenuhi
$\chi_\rho(e) = \dim V$ dan $\chi_\rho(hgh^{-1}) = \chi_\rho(g)$
(sebab trace bersifat invarian terhadap konjugasi): jadi karakter merupakan \emph{fungsi
kelas}\index{fungsi kelas} --- yaitu unsur ruang
$\mathcal{CF}(G)$ berisi fungsi yang konstan pada \omterm{ex:b3:groups:actions}{kelas konjugasi},
yang dilengkapi hasil kali dalam Hermite
\[
\langle \varphi, \psi\rangle = \frac1{\abs G}\sum_{g \in G}
\overline{\varphi(g)}\,\psi(g).
\]
\end{definition}

\begin{proposition}\label{prop:b3:representations:charbasics}
Pemetaan $\rho(g)$ dapat didiagonalkan dengan nilai eigen berupa akar satuan;
$\chi_\rho(g^{-1}) = \overline{\chi_\rho(g)}$, dan
$\abs{\chi_\rho(g)} \leq \chi_\rho(e)$ dengan kesamaan jika dan hanya jika
$\rho(g)$ berupa skalar. \omterm{def:b3:representations:character}{Karakter} menjumlah pada jumlah langsung: $\chi_{V
\oplus W} = \chi_V + \chi_W$.
\end{proposition}

\begin{proof}
Berlaku $\rho(g)^N = \mathrm{id}$ untuk $N = \abs G$ (menurut Lagrange): jadi $\rho(g)$
menghapus $X^N - 1$, yang terpecah dengan akar sederhana, sehingga ia dapat
didiagonalkan dengan nilai eigen $\lambda_i \in \mu_N$. Lalu
$\chi(g^{-1}) = \sum \lambda_i^{-1} = \sum\bar\lambda_i =
\overline{\chi(g)}$, dan $\abs{\chi(g)} = \abs{\sum\lambda_i}
\leq \dim V$, dengan kesamaan pada ketaksamaan segitiganya jika dan hanya jika semua
$\lambda_i$ sama, yakni $\rho(g) = \lambda\,\mathrm{id}$.
Keaditifannya pada jumlah langsung: lihat trace bloknya.
\end{proof}

\begin{lemma}\label{lem:b3:representations:homtrace}
Misalkan $(V, \rho)$, $(W, \sigma)$ dua \omterm{def:b3:representations:rep}{representasi}. Operator
perata-rataan pada $\operatorname{Hom}(W, V)$,
\[
c(A) = \frac1{\abs G}\sum_{g}\rho(g)\,A\,\sigma(g)^{-1},
\]
merupakan proyeksi pada $\operatorname{Hom}_G(W, V)$, dan untuk
pemetaan $\Phi_g \colon A \mapsto \rho(g)A\sigma(g)^{-1}$ berlaku
$\operatorname{tr}\Phi_g =
\chi_\rho(g)\,\overline{\chi_\sigma(g)}$.
\end{lemma}

\begin{proof}
Pemetaan $c(A)$ bersifat ekuivarian (indeks ulang jumlahnya seperti pada Maschke), dan $c$
menetapkan pemetaan ekuivarian (sebab setiap sukunya sama dengan $A$): jadi $c$ merupakan
proyeksi berpeta $\operatorname{Hom}_G(W, V)$. Tracenya: dalam
basis, $\Phi_g(A) = BAC$ dengan $B = \rho(g)$, $C = \sigma(g)^{-1}$;
pada basis matriks $(E_{kl})$ berlaku $BE_{kl}C = \sum_{m,n}
b_{mk}c_{ln}E_{mn}$, sehingga koefisien $E_{kl}$ pada
$\Phi_g(E_{kl})$ adalah $b_{kk}c_{ll}$: jadi $\operatorname{tr}\Phi_g =
\sum_{k,l}b_{kk}c_{ll} = \operatorname{tr}(B)
\operatorname{tr}(C) = \chi_\rho(g)\chi_\sigma(g^{-1})$, dan
$\chi_\sigma(g^{-1}) = \overline{\chi_\sigma(g)}$.
\end{proof}

\begin{theorem}[Relasi ortogonalitas pertama]
\label{thm:b3:representations:orthogonality}
Misalkan $\rho, \sigma$ \emph{\omterm{def:b3:representations:rep}{tak tereduksi}}. Maka\index{ortogonalitas
karakter}
\[
\langle\chi_\sigma, \chi_\rho\rangle =
\begin{cases}
1 & \text{jika } \rho \cong \sigma,\\
0 & \text{bila tidak:}
\end{cases}
\]
jadi \omterm{def:b3:representations:character}{karakter} \omterm{def:b3:representations:rep}{tak tereduksi} membentuk keluarga ortonormal di
$\mathcal{CF}(G)$.
\end{theorem}

\begin{proof}
Trace sebuah proyeksi sama dengan dimensi petanya:
\[
\dim\operatorname{Hom}_G(W, V) = \operatorname{tr} c
= \frac1{\abs G}\sum_g \operatorname{tr}\Phi_g
= \frac1{\abs G}\sum_g \chi_\rho(g)\overline{\chi_\sigma(g)}
= \langle \chi_\sigma, \chi_\rho\rangle,
\]
dan lema Schur menilai ruas kirinya sebagai $\delta_{\rho \cong
\sigma}$.
\end{proof}

\begin{corollary}\label{cor:b3:representations:multiplicity}
Uraikan $V \cong \bigoplus_i V_i^{\oplus m_i}$ atas \omterm{def:b3:representations:rep}{representasi tak tereduksi}
yang berbeda ($V_i \not\cong V_j$). Maka $m_i = \langle
\chi_{V_i}, \chi_V\rangle$: jadi multiplisitasnya --- dan karenanya
\omterm{def:b3:representations:rep}{representasinya} sampai isomorfisma --- ditentukan oleh
\omterm{def:b3:representations:character}{karakternya}. Lebih lanjut $\langle \chi_V, \chi_V\rangle = \sum_i
m_i^2$; khususnya $V$ \omterm{def:b3:representations:rep}{tak tereduksi} jika dan hanya jika $\langle\chi_V,
\chi_V\rangle = 1$.
\end{corollary}

\begin{proof}
Berlaku $\chi_V = \sum m_i\chi_{V_i}$
(\cref{prop:b3:representations:charbasics}); ambillah hasil kali dalam
dengan setiap $\chi_{V_i}$ lalu pakailah keortonormalannya. Dua
\omterm{def:b3:representations:rep}{representasi} dengan \omterm{def:b3:representations:character}{karakter} yang sama mempunyai multiplisitas yang sama,
sehingga keduanya isomorfik.
\end{proof}

\begin{theorem}[Representasi reguler]
\label{thm:b3:representations:regular}
Misalkan $\chi_1, \dots, \chi_r$ \omterm{def:b3:representations:character}{karakter} \omterm{def:b3:representations:rep}{tak tereduksi} yang
berbeda, berderajat $n_i = \chi_i(e)$. \omterm{ex:b3:representations:examples}{Representasi
reguler} terurai dengan multiplisitas $m_i = n_i$;
akibatnya
\[
\sum_{i=1}^{r} n_i^2 = \abs G,
\qquad
\sum_i n_i\chi_i(g) = 0 \quad (g \neq e).
\]
\end{theorem}

\begin{proof}
\omterm{def:b3:representations:character}{Karakter} regulernya: $\chi_{\mathrm{reg}}(g) =
\#\{h : gh = h\}$, yang bernilai $\abs G$ untuk $g = e$ dan $0$
bila tidak. Jadi $m_i = \langle \chi_i,\chi_{\mathrm{reg}}\rangle =
\frac1{\abs G}\overline{\chi_i(e)}\,\abs G = n_i$. Memasukkan nilai
$\chi_{\mathrm{reg}} = \sum n_i\chi_i$ pada $e$ dan pada $g \neq e$
memberikan kedua identitas yang ditampilkan itu.
\end{proof}

\begin{theorem}\label{thm:b3:representations:numberirr}
\omterm{def:b3:representations:character}{Karakter} \omterm{def:b3:representations:rep}{tak tereduksi} membentuk \emph{basis} ortonormal bagi
$\mathcal{CF}(G)$: jadi banyaknya \omterm{def:b3:representations:rep}{representasi tak tereduksi}
sama dengan banyaknya \omterm{ex:b3:groups:actions}{kelas konjugasi} $G$.
\end{theorem}

\begin{proof}
Tersisa kelengkapannya: misalkan $f \in \mathcal{CF}(G)$
ortogonal terhadap setiap $\chi_i$; kita tunjukkan $f = 0$. Untuk sebuah
\omterm{def:b3:representations:rep}{representasi} $(V,\rho)$, tetapkan $T_{f,\rho} = \frac{1}{\abs G}
\sum_g \overline{f(g)}\,\rho(g)$. Operator ini ekuivarian: sebab untuk $h \in
G$,
\[
\rho(h)T_{f,\rho}\rho(h)^{-1} = \frac1{\abs G}\sum_g
\overline{f(g)}\rho(hgh^{-1})
= \frac1{\abs G}\sum_{g'}\overline{f(h^{-1}g'h)}\rho(g') =
T_{f,\rho}
\]
(karena $f$ \omterm{def:b3:representations:character}{fungsi kelas}). Jika $\rho$ \omterm{def:b3:representations:rep}{tak tereduksi} berderajat
$n$, maka Schur memberikan $T_{f,\rho} = \lambda\,\mathrm{id}$ dengan
\[
\lambda = \frac{\operatorname{tr}T_{f,\rho}}{n}
= \frac{1}{n\abs G}\sum_g \overline{f(g)}\,\chi_\rho(g)
= \frac1n\,\langle f, \chi_\rho\rangle = 0 .
\]
Jadi $T_{f,\rho} = 0$ pada setiap \omterm{def:b3:representations:rep}{representasi tak tereduksi}, sehingga (lewat jumlah langsung) pada
\emph{setiap} \omterm{def:b3:representations:rep}{representasi} --- khususnya pada yang reguler.
Terapkan ia pada vektor basis $e_e$: $0 = T_{f,\mathrm{reg}}e_e =
\frac1{\abs G}\sum_g\overline{f(g)}e_g$, yang memaksa setiap
$\overline{f(g)} = 0$. Dengan demikian keluarga ortonormal
$(\chi_i)$ merentang $\mathcal{CF}(G)$, yang dimensinya sama dengan
banyaknya \omterm{ex:b3:groups:actions}{kelas konjugasi}.
\end{proof}

\begin{corollary}[Ortogonalitas kolom]
\label{cor:b3:representations:column}
Untuk $g, h \in G$:
\[
\sum_{i=1}^{r}\overline{\chi_i(g)}\,\chi_i(h) =
\begin{cases}
\abs{Z_G(g)} & \text{jika } g, h \text{ sekawan},\\
0 & \text{bila tidak.}
\end{cases}
\]
\end{corollary}

\begin{proof}
Misalkan $g_1, \dots, g_r$ mewakili kelasnya, dan $c_j$ ukuran
kelasnya. Matriks $r \times r$ dengan $U_{ij} =
\sqrt{c_j/\abs G}\;\chi_i(g_j)$ mempunyai baris yang ortonormal
(\cref{thm:b3:representations:orthogonality} yang ditulis per kelas:
$\sum_j \frac{c_j}{\abs G}\chi_i(g_j)\overline{\chi_{i'}(g_j)} =
\delta_{ii'}$), yakni $UU^* = I$; dan matriks persegi dengan $UU^* =
I$ juga memenuhi $U^*U = I$: jadi kolomnya ortonormal, dan itu terurai menjadi
identitas yang ditampilkan (sebab $\abs G/c_j = \abs{Z_G(g_j)}$, yaitu
relasi orbit--stabilisator untuk konjugasi).
\end{proof}

\begin{proposition}[Karakter berdimensi satu; pengangkatan]
\label{prop:b3:representations:onedim}
(a) Grup $G$ abelian jika dan hanya jika semua \omterm{def:b3:representations:rep}{representasi tak tereduksinya}
berderajat $1$; sedangkan banyaknya \omterm{def:b3:representations:character}{karakter} berderajat $1$ pada $G$ apa pun adalah
$[G : D(G)]$ (yaitu \omterm{def:b3:representations:character}{karakter} abelianisasinya).
(b) Jika $N \trianglelefteq G$, maka \omterm{def:b3:representations:rep}{representasi tak tereduksi} dari
$G/N$ terangkat (dengan mengomposisikannya dengan $G \to G/N$) tepat menjadi
\omterm{def:b3:representations:rep}{representasi tak tereduksi} $G$ yang kernelnya memuat $N$.
\end{proposition}

\begin{proof}
(a) Jika $G$ abelian, setiap kelasnya beranggota tunggal: sehingga $r = \abs G$,
dan $\sum n_i^2 = \abs G$ memaksa semua $n_i = 1$; sebaliknya jika
semua $n_i = 1$, maka \omterm{ex:b3:representations:examples}{representasi regulernya} berupa jumlah \omterm{def:b3:representations:rep}{representasi}
berdimensi satu, sehingga $\rho_{\mathrm{reg}}(G)$ dapat
didiagonalkan serentak, jadi komutatif, dan
$\rho_{\mathrm{reg}}$ bersifat setia: berarti $G$ abelian. \omterm{def:b3:representations:rep}{Representasi}
berderajat $1$ adalah morfisma $G \to \C^\times$ dengan sasaran abelian:
sehingga terfaktorkan melalui $G^{\mathrm{ab}} = G/D(G)$
(\cref{exo:b3:groups:9}), dan \omterm{def:b3:representations:character}{karakter} berbeda dari $G^{\mathrm{ab}}$
yang abelian berjumlah $\abs{G^{\mathrm{ab}}}$ (sebab ia mempunyai sebanyak
itu kelas, semuanya berderajat $1$). (b) Komposisi dengan
proyeksinya mempertahankan ketaktereduksian (sebab subruang invariannya
berkorespondensi), dan \omterm{def:b3:representations:rep}{representasi} yang trivial pada $N$ terfaktorkan melalui
\omterm{thm:b3:groups:quotient}{grup kuosiennya} (\cref{thm:b3:groups:firstiso}).
\end{proof}

\section{Tabel karakter}

\begin{definition}\label{def:b3:representations:table}
Yang disebut \emph{tabel karakter}\index{tabel karakter} dari $G$ adalah
matriks $r \times r$ $\bigl(\chi_i(g_j)\bigr)$: barisnya diindeks oleh
\omterm{def:b3:representations:character}{karakter} \omterm{def:b3:representations:rep}{tak tereduksi}, kolomnya oleh \omterm{ex:b3:groups:actions}{kelas konjugasi} (dengan ukurannya
ditampilkan). Barisnya ortonormal terhadap hasil kali berbobot itu, dan
kolomnya ortogonal (\cref{cor:b3:representations:column}): jadi
tabelnya sangat berlebih ditentukan, dan justru itulah yang membuatnya
terhitungkan.
\end{definition}

\begin{example}[Tabel $S_3$]\label{ex:b3:representations:s3}
Kelasnya: $e$ (berukuran 1), transposisi (3), siklus-$3$ (2); jadi ada $r
= 3$ \omterm{def:b3:representations:rep}{representasi tak tereduksi}, berderajat $n_i$ dengan $\sum n_i^2 = 6$: yakni $1, 1,
2$. Berderajat $1$: yang trivial $\mathbf 1$ dan tanda $\varepsilon$.
Baris terakhirnya menyusul dari ortogonalitas kolom (atau dari
$\chi_{\mathrm{std}} = \chi_{\mathrm{perm}} - \mathbf 1$):
\[
\begin{array}{c|ccc}
S_3 & e & (1\,2)\ [3] & (1\,2\,3)\ [2]\\
\hline
\mathbf 1 & 1 & 1 & 1\\
\varepsilon & 1 & -1 & 1\\
\chi_{\mathrm{std}} & 2 & 0 & -1
\end{array}
\]
Periksa: $\langle\chi_{\mathrm{std}},\chi_{\mathrm{std}}\rangle =
\frac{1}{6}(4 + 0 + 2) = 1$: jadi \omterm{def:b3:representations:rep}{tak tereduksi}.
\end{example}

\begin{example}[Tabel $S_4$]\label{ex:b3:representations:s4}
Kelasnya: $e$ [1], transposisi [6], transposisi ganda [3],
siklus-$3$ [8], siklus-$4$ [6]: jadi lima \omterm{def:b3:representations:rep}{representasi tak tereduksi}, dengan $\sum n_i^2 =
24$ dan dua \omterm{def:b3:representations:character}{karakter} berderajat $1$ ($\mathbf 1, \varepsilon$; sebab $[S_4 :
D(S_4)] = [S_4 : A_4] = 2$): derajatnya $1, 1, 2, 3, 3$. Yang
berderajat $2$ terangkat dari $S_4/V \cong S_3$
(\cref{prop:b3:representations:onedim}(b), dengan $V$ grup Klein);
sedangkan yang berderajat $3$: \omterm{def:b3:representations:rep}{representasi} baku dan puntirannya oleh
$\varepsilon$:
\[
\begin{array}{c|ccccc}
S_4 & e\,[1] & (1\,2)\,[6] & (1\,2)(3\,4)\,[3] &
(1\,2\,3)\,[8] & (1\,2\,3\,4)\,[6]\\
\hline
\mathbf 1 & 1 & 1 & 1 & 1 & 1\\
\varepsilon & 1 & -1 & 1 & 1 & -1\\
\chi_2 & 2 & 0 & 2 & -1 & 0\\
\chi_{\mathrm{std}} & 3 & 1 & -1 & 0 & -1\\
\varepsilon\chi_{\mathrm{std}} & 3 & -1 & -1 & 0 & 1
\end{array}
\]
($\chi_{\mathrm{std}}(g) = \operatorname{fix}(g) - 1$; sedangkan
$\chi_2$ menilai tabel $S_3$ pada peta setiap kelasnya
modulo $V$.) Semua pemeriksaan ortogonalitas baris dan kolomnya lulus ---
menjalankan dua di antaranya adalah pemanasan \cref{exo:b3:representations:3}.
\end{example}

\begin{method}\label{met:b3:representations:table}
Untuk membangun \omterm{def:b3:representations:table}{tabel karakter}: (1) daftarkan \omterm{ex:b3:groups:actions}{kelas konjugasi} beserta ukurannya;
(2) cacahlah \omterm{def:b3:representations:character}{karakter} berderajat $1$ lewat $G/D(G)$ lalu tuliskan; (3)
carilah derajat sisanya dari $\sum n_i^2 = \abs G$ (lewat kombinatorika
bilangan bulat kecil); (4) perolehlah \omterm{def:b3:representations:rep}{representasi tak tereduksi} yang murah: angkat dari
\omterm{thm:b3:groups:quotient}{grup kuosiennya}, kurangkan $\mathbf 1$ dari \omterm{def:b3:representations:character}{karakter} permutasinya
(periksa $\langle\chi,\chi\rangle = 1$), kalikan \omterm{def:b3:representations:character}{karakter} yang sudah dikenal
dengan yang berderajat $1$; (5) tuntaskan baris yang belum diketahui lewat ortogonalitas
kolom --- sebab setiap kolom ortogonal terhadap kolom yang
sudah lengkap, dan kolom $e$ membawa derajatnya. Periksalah
semuanya dengan sapuan ortogonalitas yang penuh.
\end{method}

\begin{omfigure}
\begin{tikzpicture}[scale=0.52]
  \draw[thick] (0,0) rectangle (6,6);
  \fill[omDef!25] (0,5) rectangle (1,6);
  \fill[omProp!25] (1,4) rectangle (2,5);
  \fill[omThm!20] (2,0) rectangle (6,4);
  \draw (0,5) rectangle (1,6);
  \draw (1,4) rectangle (2,5);
  \draw (2,0) rectangle (6,4);
  \node[font=\small] at (0.5,5.5) {$1$};
  \node[font=\small] at (1.5,4.5) {$1$};
  \node[font=\small] at (4,2) {$M_2(\C)$};
\end{tikzpicture}

{\small \omterm{ex:b3:representations:examples}{Representasi reguler} $S_3$, setelah didiagonalkan atas blok:
$\C[S_3] \cong \C \times \C \times M_2(\C)$, dengan dimensi $1 + 1 +
4 = 6 = \abs{S_3}$. Secara umum $\C[G] \cong \prod_i
M_{n_i}(\C)$: jadi identitas $\sum n_i^2 = \abs G$ merupakan pernyataan
tentang blok matriks.}
\end{omfigure}

\section{Latihan}

\begin{exercise}[$\star$]\label{exo:b3:representations:1}
(a) Tunjukkan bahwa \omterm{def:b3:representations:character}{karakter} \omterm{def:b3:representations:rep}{tak tereduksi} dari $\Z/n\Z$ adalah
$\chi_k(\bar m) = \eu^{2\iu\pi km/n}$, $k = 0, \dots, n-1$, lalu
tuliskan \omterm{def:b3:representations:table}{tabel karakter} $\Z/4\Z$.
(b) Periksalah kedua relasi ortogonalitasnya pada tabel itu --- lalu kenalilah
matriksnya: di manakah buku ini pernah menjumpainya?
\end{exercise}

\begin{exercise}[$\star$]\label{exo:b3:representations:2}
Misalkan $G$ beraksi pada himpunan berhingga $X$ dan $\chi$ \omterm{def:b3:representations:character}{karakter}
\omterm{def:b3:representations:rep}{representasi} permutasi $\C^X$.
(a) Tunjukkan $\chi(g) = \abs{\operatorname{Fix}_X(g)}$ dan
$\langle \mathbf 1, \chi\rangle = \#\{\text{orbit}\}$ ---
jadi lema pencacahan Burnside (\cref{exo:b3:groups:5}) merupakan perhitungan
\omterm{def:b3:representations:character}{karakter}.
(b) Andaikan aksinya transitif, sehingga $\chi = \mathbf 1 +
\psi$. Tunjukkan bahwa $\psi$ \omterm{def:b3:representations:rep}{tak tereduksi} jika dan hanya jika aksinya
\emph{$2$-transitif} (yakni transitif pada pasangan terurut titik yang
berbeda). \emph{(Hitunglah $\langle\chi,\chi\rangle$ sebagai banyaknya
\omterm{def:b3:groups:action}{orbit} pada $X \times X$.)}
(c) Simpulkan bahwa \omterm{def:b3:representations:rep}{representasi} baku $S_n$ (untuk $n \geq
2$) bersifat \omterm{def:b3:representations:rep}{tak tereduksi}.
\end{exercise}

\begin{exercise}[$\star$]\label{exo:b3:representations:3}
Bangun ulang tabel $S_3$ dari nol dengan mengikuti
\cref{met:b3:representations:table}, lalu periksalah dua relasi ortogonalitas
baris dan dua relasi ortogonalitas kolom pada tabel $S_4$ di
\cref{ex:b3:representations:s4}. Uraikan \omterm{def:b3:representations:character}{karakter} permutasi
$S_4$ yang beraksi pada $\{1,2,3,4\}$ dan \omterm{def:b3:representations:character}{karakter}
$\chi_{\mathrm{std}}^2$ (kuadrat titik demi titik) atas semua \omterm{def:b3:representations:rep}{representasi tak tereduksi}.
\end{exercise}

\begin{exercise}[$\star\star$]\label{exo:b3:representations:4}
Hitunglah \omterm{def:b3:representations:table}{tabel karakter} $D_4$ dan $Q_8$. Simpulkan
bahwa dua grup yang tak isomorfik dapat mempunyai \omterm{def:b3:representations:table}{tabel karakter} yang
identik --- data teori grup apa yang meski demikian masih ditangkap
tabel itu pada pasangan ini (orde pusatnya, abelianisasinya,
banyaknya involusinya)? Manakah di antaranya yang \emph{gagal}
ditangkap?
\end{exercise}

\begin{exercise}[$\star\star$]\label{exo:b3:representations:5}
\omterm{def:b3:representations:table}{Tabel karakter} $A_4$: kelasnya $e$ [1], transposisi ganda
[3], dan \emph{dua} kelas siklus-$3$ [4], [4].
(a) Jelaskan terpecahnya siklus-$3$ itu (bandingkan \omterm{ex:b3:groups:actions}{sentralisatornya}
di $S_4$ dan di $A_4$, seperti pada \cref{exo:b3:groups:11}).
(b) Carilah ketiga \omterm{def:b3:representations:character}{karakter} berderajat $1$ (lewat $A_4/V \cong
\Z/3\Z$) dan \omterm{def:b3:representations:character}{karakter} berderajat $3$ (dengan membatasi
$\chi_{\mathrm{std}}$ dari $S_4$), lalu susunlah tabelnya.
(c) Bacalah \omterm{def:b3:groups:normal}{subgrup normal} $A_4$ dari tabel itu
(lewat kernel $\{g : \chi(g) = \chi(e)\}$ dan irisannya).
\end{exercise}

\begin{exercise}[$\star\star$]\label{exo:b3:representations:6}
(a) Tunjukkan bahwa $\ker\chi = \{g : \chi(g) = \chi(e)\}$ merupakan
kernel \omterm{def:b3:representations:rep}{representasi} yang mendasarinya
(\cref{prop:b3:representations:charbasics}, pada kasus kesamaannya).
(b) Tunjukkan bahwa setiap \omterm{def:b3:groups:normal}{subgrup normal} $G$ merupakan irisan
kernel \omterm{def:b3:representations:character}{karakter} \omterm{def:b3:representations:rep}{tak tereduksi}. \emph{(Representasikan $G/N$
secara setia: pakai \omterm{ex:b3:representations:examples}{representasi regulernya}.)}
(c) Simpulkan: $G$ sederhana jika dan hanya jika $\ker\chi_i = \{e\}$ untuk setiap
$\chi_i$ \omterm{def:b3:representations:rep}{tak tereduksi} yang tak trivial --- jadi kesederhanaan terbaca dari
\omterm{def:b3:representations:table}{tabel karakternya}.
\end{exercise}

\begin{exercise}[$\star\star$]\label{exo:b3:representations:7}
Derajat untuk $\abs G = 8$: tunjukkan bahwa grup tak abelian berorde
$8$ berpola derajat $(1,1,1,1,2)$, dan bahwa \omterm{def:b3:representations:rep}{representasinya} yang
berderajat $2$ bersifat setia. Secara lebih umum tunjukkan bahwa grup
tak abelian berorde $p^3$ berpola $(1^{\,p^2},
p, \dots, p)$ dengan $p^2$ angka satu dan $p - 1$ \omterm{def:b3:representations:character}{karakter} berderajat
$p$. \emph{(Pakailah $[G : D(G)]$ dan $\sum n_i^2 = \abs G$; di sini
$D(G) = Z(G)$ berorde $p$.)}
\end{exercise}

\begin{exercise}[$\star\star\star$]\label{exo:b3:representations:8}
Untuk grup berhingga $G, H$: tunjukkan bahwa \omterm{def:b3:representations:character}{fungsi kelas}
$\chi(g)\psi(h)$ pada $G \times H$, untuk $\chi, \psi$ berupa \omterm{def:b3:representations:character}{karakter}
\omterm{def:b3:representations:rep}{tak tereduksi} dari $G, H$, tepat merupakan \omterm{def:b3:representations:character}{karakter} \omterm{def:b3:representations:rep}{tak tereduksi}
$G \times H$. \emph{(Keortonormalannya berupa perhitungan langsung;
kelengkapannya lewat pencacahan kelas.)} Simpulkan \omterm{def:b3:representations:table}{tabel karakter}
$\Z/2\Z \times \Z/2\Z$ lalu turunkan kembali
\cref{prop:b3:representations:onedim}(a) untuk grup abelian
berhingga lewat teorema struktur.
\end{exercise}

\begin{exercise}[$\star\star\star$]\label{exo:b3:representations:9}
\omterm{def:b3:representations:table}{Tabel karakter} $A_5$ (dengan kelas berukuran $1, 15, 20, 12,
12$ dari \cref{exo:b3:groups:11}):
(a) Tunjukkan bahwa derajatnya $1, 3, 3, 4, 5$ \emph{(satu-satunya
penyelesaian $\sum n_i^2 = 60$ dengan $n_1 = 1$ dan, memakai
\cref{exo:b3:representations:6}(c) beserta kesederhanaannya, tanpa
$n_i = 1$ yang lain)}.
(b) Bangunlah \omterm{def:b3:representations:character}{karakter} berderajat $4$ (lewat aksi permutasi pada
$5$ titik) dan \omterm{def:b3:representations:character}{karakter} berderajat $5$ (aksi pada keenam
subgrup Sylow-$5$ memberikan derajat $6 = 1 + 5$; periksa
ketaktereduksiannya), lalu lengkapi kedua baris berderajat $3$ lewat ortogonalitas
kolom: entri rasio emas $\frac{1\pm\sqrt5}2$ akan muncul
pada kelas siklus-$5$.
(c) Periksalah dari tabel yang sudah jadi bahwa $A_5$ sederhana
(\cref{exo:b3:representations:6}(c)).
\end{exercise}

\begin{exercise}[$\star\star$]\label{exo:b3:representations:10}
Misalkan $\rho$ \omterm{def:b3:representations:rep}{representasi tak tereduksi} berderajat $n$ dan $z
\in Z(G)$. Tunjukkan $\rho(z) = \lambda_z\,\mathrm{id}$ dengan
$\lambda \colon Z(G) \to \C^\times$ sebuah morfisma (yaitu
\emph{\omterm{def:b3:representations:character}{karakter} pusat}), lalu simpulkan $\abs{\chi(z)} = n$ untuk
$z$ yang berada di pusat. Terapannya: jika $G$ mempunyai \omterm{def:b3:representations:rep}{representasi tak tereduksi}
yang setia, maka $Z(G)$ bersifat siklik.
\end{exercise}

\begin{exercise}[$\star\star$]\label{exo:b3:representations:11}
(Proyeksi isotipik) Misalkan $(V, \rho)$ sebuah \omterm{def:b3:representations:rep}{representasi}
$G$ dan $\chi_i$ \omterm{def:b3:representations:character}{karakter} \omterm{def:b3:representations:rep}{tak tereduksi} berderajat $n_i$.
Definisikan
\[
p_i = \frac{n_i}{\abs G}\sum_{g\in G}
\overline{\chi_i(g)}\,\rho(g) \;\in\; \mathcal L(V).
\]
(a) Tunjukkan bahwa $p_i$ bersifat ekuivarian-$G$, lalu hitunglah
pembatasannya pada subrepresentasi \omterm{def:b3:representations:rep}{tak tereduksi} $W \subseteq
V$ yang berkarakter $\chi_j$: hasilnya $\delta_{ij}\,
\mathrm{id}_W$ \emph{(lewat Schur; ambillah trace untuk mengenali
skalarnya)}.
(b) Simpulkan bahwa $p_i$ merupakan proyeksi pada jumlah $V_i$ dari
semua subrepresentasi \omterm{def:b3:representations:rep}{tak tereduksi} berkarakter $\chi_i$ (yaitu
\emph{komponen isotipiknya}), bahwa $\sum_ip_i =
\mathrm{id}_V$, dan bahwa penguraian $V =
\bigoplus_iV_i$ bersifat kanonik --- tidak seperti pemecahan yang lebih halus
dari setiap $V_i$ atas \omterm{def:b3:representations:rep}{representasi tak tereduksi}.
(c) Untuk \omterm{ex:b3:representations:examples}{representasi reguler} $S_3$ dan \omterm{def:b3:representations:character}{karakter} tanda
$\varepsilon$, tuliskan $p_\varepsilon$ secara eksplisit
sebagai \omterm{def:b3:galois:algebraic}{unsur aljabar} grupnya lalu periksalah $p_\varepsilon^2
= p_\varepsilon$ dengan tangan.
\end{exercise}

\begin{exercise}[$\star\star$]\label{exo:b3:representations:12}
(Membaca sebuah tabel) \omterm{def:b3:representations:table}{Tabel karakter} suatu grup $G$
berorde $24$ diketahui sebagian: ia mempunyai $5$ kelas berukuran
$1, 6, 8, 6, 3$, dan derajat $1, 1, 2, 3, 3$.
(a) Pulihkan tabel selengkapnya: kedua \omterm{def:b3:representations:character}{karakter} linearnya
(satu trivial; yang lain bernilai $-1$ tepat pada
kelas berukuran $6$ dan $6$), lalu \omterm{def:b3:representations:character}{karakter} berderajat $2$
lewat ortogonalitas kolom dengan kolom identitasnya, lalu
kedua \omterm{def:b3:representations:character}{karakter} berderajat $3$ dengan cara serupa (salah satunya $\chi_2\chi_4$).
(b) Kenali $G$ (yaitu $\cong S_4$: bandingkan kelasnya dengan tipe
siklusnya), lalu ambillah dari tabel itu \omterm{def:b3:groups:normal}{subgrup normalnya} lewat
\cref{exo:b3:representations:6}: yakni kernel $\chi_2$ (berindeks
$2$: $A_4$) dan kernel $\chi_3$ (grup Klein $V_4$), dan tak ada
lagi selain $\{e\}, G$.
(c) Jelaskan bagaimana tabel itu menunjukkan $G/V_4 \cong S_3$
\emph{(\omterm{def:b3:representations:character}{karakter} mana yang terfaktorkan melalui \omterm{thm:b3:groups:quotient}{grup kuosiennya}?)}.
\end{exercise}

\section{Soal: teorema \texorpdfstring{$p^aq^b$}{paqb}
Burnside}

\begin{problem}[{Soal akhir pekan --- keterselesaian grup
berorde $p^aq^b$}]\label{pb:b3:representations:1}
Burnside membuktikan pada 1904 bahwa setiap grup yang ordenya berfaktor prima
paling banyak dua bersifat \omterm{def:b3:groups:derived}{terselesaikan} --- sebuah pernyataan tentang grup
abstrak yang selama setengah abad hanya diketahui buktinya lewat
teori \omterm{def:b3:representations:character}{karakter}. Soal ini membangun buktinya secara penuh, dengan
merakit \cref{ch:b3:groups} (\omterm{def:b3:groups:derived}{keterselesaian}),
\cref{ch:b3:modules} (modul-$\Z$ yang \omterm{def:b3:modules:free}{dibangun secara berhingga}) dan bab
ini. Sepanjang soal, $\chi_1, \dots, \chi_r$ adalah \omterm{def:b3:representations:character}{karakter} \omterm{def:b3:representations:rep}{tak tereduksi}
$G$, dan $n_i = \chi_i(e)$.

\medskip
\noindent\textbf{Bagian I --- Bilangan bulat aljabar.} Sebuah
\emph{bilangan bulat aljabar} adalah akar sebuah polinomial \emph{monik}
di $\Z[X]$.
\begin{enumerate}
  \item Tunjukkan bahwa $\alpha$ merupakan bilangan bulat aljabar jika dan hanya jika ring
        $\Z[\alpha]$ merupakan modul-$\Z$ yang \omterm{def:b3:modules:free}{dibangun secara berhingga}.
  \item Simpulkan bahwa bilangan bulat aljabar membentuk subring
        $\C$. \emph{(Jika $\Z[\alpha], \Z[\beta]$ \omterm{def:b3:modules:free}{dibangun secara
        berhingga}, maka demikian pula $\Z[\alpha, \beta]$, dan submodul dari
        modul-$\Z$ yang \omterm{def:b3:modules:free}{dibangun secara berhingga} juga \omterm{def:b3:modules:free}{dibangun secara berhingga},
        menurut \cref{thm:b3:modules:submodule} beserta argumen penyajian
        --- atau secara langsung: submodul $\Z^n$ bersifat bebas
        berrank $\leq n$.)}
  \item Tunjukkan bahwa bilangan bulat aljabar yang \emph{rasional} pastilah
        bilangan bulat. \emph{(Lewat teorema akar rasional.)}
  \item Tunjukkan bahwa setiap nilai \omterm{def:b3:representations:character}{karakter} $\chi(g)$ merupakan
        bilangan bulat aljabar.
\end{enumerate}

\medskip
\noindent\textbf{Bagian II --- Relasi jumlah kelas.} Tetapkan sebuah
\omterm{def:b3:representations:rep}{representasi tak tereduksi} $(V, \rho)$ berderajat $n$ dan berkarakter $\chi$. Untuk sebuah
\omterm{ex:b3:groups:actions}{kelas konjugasi} $C$, misalkan $S_C = \sum_{g \in C}\rho(g) \in
\mathcal L(V)$.
\begin{enumerate}[resume]
  \item Tunjukkan bahwa $S_C$ bersifat ekuivarian, sehingga $S_C =
        \omega_C\,\mathrm{id}$ dengan
        \[
        \omega_C = \frac{\abs C\,\chi(g_C)}{n}
        \qquad (g_C \in C \text{ sembarang wakilnya}).
        \]
  \item Tunjukkan bahwa $S_CS_{C'} = \sum_{C''}
        a_{CC'C''}\,S_{C''}$ dengan $a_{CC'C''} \in \N$ mencacah,
        untuk $z \in C''$ yang tetap, pasangan $(x, y) \in C \times
        C'$ dengan $xy = z$. Simpulkan bahwa $\omega_C$ memenuhi
        $\omega_C\,\omega_{C'} = \sum_{C''}
        a_{CC'C''}\,\omega_{C''}$.
  \item Simpulkan bahwa setiap $\omega_C$ merupakan bilangan bulat aljabar.
        \emph{(Modul-$\Z$ yang dibangun oleh $1$ dan
        $\omega_C$ merupakan ring yang \omterm{def:b3:modules:free}{dibangun secara berhingga}; terapkan kriteria
        pertanyaan 1 --- lebih tepatnya tunjukkan bahwa $M = \Z\text{-rentang}
        (1, (\omega_C)_C)$ memenuhi $\omega_{C_0} M \subseteq
        M$ lalu pakailah muslihat determinan/Cayley--Hamilton, atau
        argumen subring pada pertanyaan 2.)}
  \item Simpulkan \emph{keterbagian \omterm{thm:b3:galois:finitefields}{Frobenius}}: bahwa $n_i$ membagi
        $\abs G$ untuk setiap derajat \omterm{def:b3:representations:rep}{tak tereduksi}.
        \emph{(Hitunglah $\frac{\abs G}{n} = \frac{\abs G}{n}
        \langle\chi,\chi\rangle = \sum_C \omega_C\,
        \overline{\chi(g_C)}$: sebuah bilangan bulat aljabar yang
        rasional.)}
\end{enumerate}

\medskip
\noindent\textbf{Bagian III --- Kriteria kesederhanaan Burnside.}
\begin{enumerate}[resume]
  \item Misalkan $\chi$ \omterm{def:b3:representations:rep}{tak tereduksi} berderajat $n$ dan $C$ sebuah kelas
        dengan $\gcd(\abs C, n) = 1$. Dengan memakai Bézout dan pertanyaan
        4--7, tunjukkan bahwa $\frac{\chi(g_C)}{n}$ merupakan bilangan bulat
        aljabar.
  \item Andaikan lebih lanjut $0 < \abs{\chi(g_C)} < n$. Tunjukkan bahwa ini
        mustahil: bilangan bulat aljabar $\alpha =
        \chi(g_C)/n$ mempunyai seluruh \emph{konjugatnya} --- yakni
        bilangan $\alpha_\sigma = \frac1n\sigma(\chi(g_C))$ untuk
        $\sigma \in \operatorname{Gal}(\Q(\zeta_{\abs G})/\Q)$,
        yang masing-masing rata-rata $n$ akar satuan --- bermodulus
        $\leq 1$,
        sehingga hasil kali $N = \prod_\sigma\alpha_\sigma$ menjadi
        bilangan bulat aljabar rasional dengan $0 < \abs N < 1$ ---
        berilah alasan bagi setiap pernyataan itu, sambil mengutip
        \cref{thm:b3:galois:cyclotomicirred} untuk grup Galoisnya
        dan kenyataan bahwa $\sigma$ mempermutasikan akar
        satuan. Simpulkan: entah $\chi(g_C) = 0$ entah $\rho(g_C)$
        berupa skalar (\cref{prop:b3:representations:charbasics}).
  \item (Kriteria Burnside) Misalkan $C \neq \{e\}$ sebuah
        \omterm{ex:b3:groups:actions}{kelas konjugasi} yang berukuran \emph{pangkat prima} $p^k > 1$,
        dan andaikan $G$ sederhana serta tak abelian. Ortogonalitas
        kolom pada kolom $C$ terhadap kolom
        $e$ memberikan $1 + \sum_{i \geq 2} n_i\chi_i(g_C) = 0$. Tunjukkan
        bahwa ada $\chi_i$ tak trivial dengan $p \nmid n_i$ yang
        $\chi_i(g_C) \neq 0$ \emph{(sebab kalau tidak, $\frac1p$ akan menjadi
        bilangan bulat aljabar)}; lalu menurut pertanyaan 10, $\rho_i(g_C)$
        berupa skalar; turunkanlah kontradiksi dengan kesederhanaannya
        \emph{(himpunan $g$ dengan $\rho_i(g)$ berupa skalar merupakan
        \omterm{def:b3:groups:normal}{subgrup normal}; pakailah kesetiaannya dari
        \cref{exo:b3:representations:6})}. Simpulkan: \emph{tak ada
        \omterm{def:b3:groups:simple}{grup sederhana} tak abelian yang mempunyai \omterm{ex:b3:groups:actions}{kelas konjugasi} berukuran pangkat
        prima $> 1$.}
\end{enumerate}

\medskip
\noindent\textbf{Bagian IV --- Teoremanya.}
\begin{enumerate}[resume]
  \item Misalkan $\abs G = p^aq^b$ dengan $a + b \geq 1$. Jika $G$
        sederhana, tunjukkan bahwa ia abelian: ambil $z \neq e$ di
        pusat sebuah subgrup Sylow-$q$
        (\cref{thm:b3:groups:pfixed}) lalu tinjau ukuran
        \omterm{ex:b3:groups:actions}{kelas konjugasinya} $[G : Z_G(z)]$, yang berupa pangkat $p$
        (mengapa?); lalu terapkan pertanyaan 11.
  \item Simpulkan dengan induksi pada $\abs G$: \emph{setiap grup
        berorde $p^aq^b$ bersifat \omterm{def:b3:groups:derived}{terselesaikan}} (Burnside). Mengapa
        argumennya runtuh untuk tiga bilangan prima --- dan mengapa memang harus runtuh, mengingat
        $\abs{A_5} = 2^2\cdot3\cdot5$?
\end{enumerate}

\medskip
\noindent\textbf{Bagian V --- \omterm{def:b3:representations:table}{Tabel karakter} $A_5$.} Grup
terkecil yang tak tersentuh teorema Burnside layak memperoleh potret
lengkapnya; semua yang di bawah hanya memakai bab ini ditambah
\cref{exo:b3:groups:11}.
\begin{enumerate}[resume]
  \item Ingat kembali dari \cref{exo:b3:groups:11} kelima \omterm{ex:b3:groups:actions}{kelas konjugasi}
        $A_5$: $\{e\}$, $15$ transposisi
        ganda, $20$ siklus tiga, dan dua kelas
        berisi $12$ siklus lima masing-masing, yang diwakili $c =
        (1\,2\,3\,4\,5)$ dan $c^2$. Jelaskan mengapa siklus lima itu
        terpecah menjadi dua kelas-$A_5$ padahal di dalam $S_5$
        keduanya membentuk satu kelas saja.
  \item Tunjukkan bahwa derajat \omterm{def:b3:representations:rep}{tak tereduksi} $A_5$ tepat berupa
        $1, 3, 3, 4, 5$: pakailah $\sum_in_i^2 = 60$ dengan $r = 5$
        kelas, beserta kenyataan bahwa $A_5$ bersifat \omterm{prop:b3:galois:perfect}{sempurna} ($D(A_5) =
        A_5$), sehingga \omterm{def:b3:representations:character}{karakter} trivial menjadi satu-satunya \omterm{def:b3:representations:character}{karakter} linearnya;
        lalu singkirkan setiap multihimpunan lainnya \emph{(tulislah $59$
        sebagai jumlah empat kuadrat bilangan bulat $\geq 2$: periksalah bahwa
        hal itu terjadi hanya dengan satu cara yang semua sukunya masuk akal
        sebagai derajat)}.
  \item Misalkan $\pi$ \omterm{def:b3:representations:character}{karakter} permutasi $A_5$ pada
        $\{1, \dots, 5\}$: $\pi(g) = \#\operatorname{Fix}(g)$,
        dengan nilai $5, 1, 2, 0, 0$ pada kelima kelasnya. Hitunglah
        $\langle\pi, \mathbf 1\rangle$ dan $\langle\pi,
        \pi\rangle$, lalu simpulkan bahwa $\chi_4 = \pi - \mathbf 1$
        \omterm{def:b3:representations:rep}{tak tereduksi} berderajat $4$, dengan nilai $4, 0, 1, -1,
        -1$.
  \item Permainan yang sama pada $10$ pasangan tak terurut $\{i, j\}$: di sini
        cacahan titik tetapnya $10, 2, 1, 0, 0$. Hitunglah
        $\langle\pi_{10}, \pi_{10}\rangle$, $\langle\pi_{10},
        \mathbf 1\rangle$ dan $\langle\pi_{10}, \chi_4\rangle$,
        simpulkan penguraian $\pi_{10} = \mathbf 1 + \chi_4
        + \chi_5$, lalu perolehlah $\chi_5$ yang \omterm{def:b3:representations:rep}{tak tereduksi} berderajat
        $5$ dengan nilai $5, 1, -1, 0, 0$.
  \item Kedua \omterm{def:b3:representations:rep}{representasi tak tereduksi} yang tersisa, $\chi_2, \chi_3$, berderajat
        $3$. Ortogonalitas kolom (setiap kolom bukan identitas terhadap kolom
        identitasnya, dan setiap kolom dengan dirinya sendiri) menentukan
        nilainya di luar siklus lima:
        tunjukkan $\chi_2(g) = \chi_3(g) = -1$ pada transposisi
        ganda dan $0$ pada siklus tiga.
  \item Pada kelas siklus lima, tetapkan $x = \chi_2(c)$ dan $y =
        \chi_2(c^2)$; simetrinya memungkinkan kita mengambil $\chi_3(c) = y$,
        $\chi_3(c^2) = x$. Dari kolom $c$ yang dipasangkan dengan
        kolom identitas dan dengan kolom $c^2$, turunkanlah
        $x + y = 1$ dan $xy = -1$ (lalu periksalah nilai $x^2 +
        y^2 = 3$ yang diberikan kolom $c$ dengan dirinya sendiri), sehingga
        \[
        \{x, y\} = \Bigl\{\frac{1 + \sqrt5}2,\ \frac{1 -
        \sqrt5}2\Bigr\} :
        \]
        yaitu rasio emas dan sekawannya. Susunlah
        \omterm{def:b3:representations:table}{tabel karakter} $A_5$ selengkapnya.
  \item Jalankan pemeriksaannya: norma baris $\chi_2$ bernilai $1$ (pakailah
        $\varphi^2 + \bar\varphi^2 = 3$), $\langle\chi_2,
        \chi_3\rangle = 0$, dan keterbagian \omterm{thm:b3:galois:finitefields}{Frobenius} (pertanyaan
        8) untuk kelima derajatnya. Di bagian mana tabel itu Anda
        \emph{melihat} perbedaan dengan $S_5$, yang seluruh
        nilai \omterm{def:b3:representations:character}{karakternya} berupa bilangan bulat rasional?
  \item Simpulkan dari tabel itu saja bahwa $A_5$ sederhana: sebuah
        \omterm{def:b3:groups:normal}{subgrup normal} berupa gabungan \omterm{ex:b3:groups:actions}{kelas konjugasi}
        yang memuat $e$ dan kardinalitasnya membagi $60$ --- periksalah
        bahwa tak ada subjumlah sejati dari $1 + 15 + 20 + 12 + 12$
        yang memuat suku $1$ dan membagi $60$. Periksa silang dengan
        kriteria pertanyaan 11: pastikan bahwa tak ada kelas $A_5$
        yang berukuran pangkat prima $> 1$.
  \item (Penutup ikosahedral) Grup $A_5$ adalah grup rotasi
        ikosahedron, dan \omterm{def:b3:representations:rep}{representasi} berderajat $3$ itu merupakan
        kedua aksi geometrinya pada $\R^3$. Periksalah identitas
        tracenya: rotasi sebesar sudut $\theta$ bertrace $1 +
        2\cos\theta$, dan $1 + 2\cos\frac{2\pi}5 =
        \frac{1+\sqrt5}2 = \varphi$. Jelaskan tanpa
        perhitungan apa pun mengapa \omterm{def:b3:representations:character}{karakter} berderajat $3$ yang lain harus
        membawa nilai sekawannya: \omterm{def:b3:galois:galois}{grup Galois}
        $\Q(\sqrt5)/\Q$ beraksi pada seluruh \omterm{def:b3:representations:table}{tabel karakternya}
        (entri demi entri), sambil mempermutasikan semua \omterm{def:b3:representations:character}{karakter} \omterm{def:b3:representations:rep}{tak tereduksinya}.
\end{enumerate}

\medskip
\noindent\textbf{Bagian VI --- Pelengkap: sebuah batas pusat dan
kuadrat tensornya.}
\begin{enumerate}[resume]
  \item (Lebih tajam daripada keterbagian) Misalkan $\chi$ \omterm{def:b3:representations:rep}{tak tereduksi}
        berderajat $n$. Tunjukkan bahwa $\abs{\chi(z)} = n$ untuk setiap
        $z \in Z(G)$ \emph{(lewat lema Schur: $\rho(z)$ berupa
        skalar berorde berhingga)}, lalu simpulkan dari
        $\langle\chi, \chi\rangle = 1$ batas
        \[
        n^2 \leq [G : Z(G)] .
        \]
        Tunjukkan bahwa grup tak abelian berorde $8$ mempunyai
        derajat \omterm{def:b3:representations:rep}{tak tereduksi} $1, 1, 1, 1, 2$ \emph{(ada lima
        \omterm{ex:b3:groups:actions}{kelas konjugasi}; tulislah $8$ sebagai jumlah lima kuadrat)}
        dan mencapai kesamaan $4 = [G : Z(G)]$; lalu periksalah batas itu
        pada $A_5$, yang pusatnya trivial.
  \item (Kuadrat tensor $\chi_2$) Untuk $g$ berorde berhingga,
        $\rho(g)$ dapat didiagonalkan dengan nilai eigen berupa akar satuan;
        simpulkan rumus \omterm{def:b3:representations:character}{karakternya}
        \[
        \chi_{\operatorname{Sym}^2 V}(g)
        = \frac{\chi(g)^2 + \chi(g^2)}2,
        \qquad
        \chi_{\Lambda^2 V}(g)
        = \frac{\chi(g)^2 - \chi(g^2)}2 .
        \]
        Terapkan rumus itu pada $\chi_2$ dari $A_5$ \emph{(perhatikan bahwa $g^2$ menjelajahi
        kelas $c^2$ ketika $g$ menjelajahi kelas
        $c$, dan sebaliknya)}: tunjukkan $\Lambda^2\chi_2 =
        \chi_2$ dan $\operatorname{Sym}^2\chi_2 = \mathbf 1 +
        \chi_5$, sehingga
        \[
        \chi_2\otimes\chi_2 = \mathbf 1 + \chi_2 + \chi_5 .
        \]
        Tafsirkan $\Lambda^2\chi_2 = \chi_2$ secara geometri lewat
        hasil kali silang pada $\R^3$.
  \item (Audit akhir atas tabelnya) Periksalah secara numerik: ortogonalitas
        kolom antara kedua kolom siklus lima
        (yakni $\varphi\bar\varphi = -1$), nilai $5 =
        \abs{Z_{A_5}(c)}$ untuk kolom $c$ terhadap dirinya sendiri,
        dan lenyapnya \omterm{def:b3:representations:character}{karakter} reguler $\sum_i
        n_i\chi_i(g) = 0$ pada masing-masing dari keempat kolom
        bukan identitas pada tabel itu.

\end{enumerate}
\end{problem}
